\section{Discussion}
\label{sec:discussion}

\paragraph{Interpretation.}
The results support a specific reading of \emph{how} an agent harness improves
instructional-design quality. AlignmentGraph-ISD leads all six baselines on the
benchmark (\cref{sec:rq1}), and the ablation attributes that lead overwhelmingly to
one mechanism: multi-agent decomposition over the alignment graph
($+3.90$ composite, $+3.29$ ADDIE, $p_{\mathrm{Holm}}=2.3\times10^{-13}$).
Decomposing generation among specialists with bounded write scope, over a shared
graph that forces each alignment edge to be created as part of generation, is what
turns constructive alignment from prompted prose into an enforced property. This is
consistent with the harness-over-model thesis~\cite{cho2026harness}: the advantage
comes from the execution environment, not from asking the model to try harder.

\paragraph{Scientific and practical implications.}
Representing alignment as first-class graph state is not education-specific: any
setting that generates many interdependent artifacts with a required correspondence
among them---technical documentation, process definitions, curricula---faces the
same invisible-defect problem and can reuse the pattern ``typed dependency graph +
bounded-scope specialists + owned verification.'' Practically, because the
advantage is structural rather than capability-driven, it is exactly the kind of
advantage expected to matter most as deployment shifts to cost-efficient
SLMs~\cite{belcak2025slm}, and the SLM run confirms this directly: the decomposition
contribution roughly doubles ($+3.9 \rightarrow +7.9$) while the harness's absolute
score barely moves and the baselines lose several points---the harness converts a
structural design into robustness against model scale-down.

\paragraph{Failure analysis: verification/repair does not move the judged score.}
Our most instructive negative result is that per-agent verification/repair
contributes essentially nothing to the composite score ($+0.26$, n.s.), even though
it demonstrably fires and repairs real pedagogical violations (\cref{sec:rq2}). Two
observations explain this and guard against over-claiming. First, on a capable
model the generated artifacts are already mostly well-aligned by construction, so
the quality-control tools have little left to fix---their activation rate is low.
Second, and more fundamentally, there is a \emph{construct gap}: the checks enforce
formal pedagogical criteria (cognitive-level match, observable rubrics,
active-practice coverage) that the ADDIE rubric, as scored by LLM judges, does not
directly reward; and an earlier, over-strict measurability check that rewrote
well-formed objective statements actively \emph{lowered} the judged score before we
narrowed it to genuinely non-observable objectives and made it audit-only. The
honest conclusion is that verification/repair is a \emph{guarantee} mechanism whose
value is visible in the residual-defect audit rather than in the judge score; on a
strong model it is close to a no-op on quality, and its case rests on the SLM
setting and on deployments where undetected alignment defects carry real cost. We
also observe a persistent trajectory-score pattern: the harness's structured
multi-agent trace is scored no higher, sometimes lower, than a clean single ReAct
loop, reflecting that trajectory judging rewards trace simplicity rather than the
structural guarantees the trace encodes.
